\begin{abstract}
A carrier--strain-gradient coupling is retained in an n-type ZnO
piezoelectric-semiconductor plate having 6\textit{mm} symmetry. Small-strain
kinematics, strain-gradient elasticity, quasistatic electrostatics, mobile-carrier
drift--diffusion, guided-wave propagation, and plate-edge reflection are retained.
The only additional constitutive interaction is
\(\Delta H=\Delta n\,H_{jkl}\varepsilon_{jk,l}\), with reciprocal terms
\(\Delta\tau_{jkl}=H_{jkl}\Delta n\) and
\(\Delta\widetilde{\mu}=H_{jkl}\varepsilon_{jk,l}\). The resulting three-displacement
calculation is reported as a 2.5D Fourier-reduced geometry with a specified
propagation-plane tangential component; it is not a strict sagittal plane-strain
calculation and no full arbitrary three-dimensional edge solution is claimed.
A Legendre spectral thickness solver and an independent edge-state formulation are
used. At the primary operating point the edge residual is
\(3.00\times10^{-15}\), the relative control-volume power balance is
\(7.31\times10^{-13}\), and the isolated backward-channel self-power at \SI{2}{\micro\metre}
is \SI{2.315}{\percent} of the incident self-power for \(H=0\). Increasing
\(|H|\) changes this isolated self-power monotonically from
\SI{2.286}{\percent} at \(H/H^*=-0.05\) to \SI{2.266}{\percent} at
\(H/H^*=+0.05\). Open-boundary runs remain singular in all tested cases and are
reported as failures, not comparisons. The synthetic \(A\), \(f\), and \(H\)
coefficients and the Qu-equivalence question are explicitly unresolved.
\end{abstract}
